% Pursuing Stacks (cotes 134-2 to 134-8): archive cote 134-2, site page 41 (PDF page 42).
% appendix (letters to Breen). Draft compared: ps1.tex:1203-1205.
% AI-transcribed by Claude Opus 5.5 (claude-opus-5-5, adaptive thinking, effort high), Message Batches API, batch ps1 (msgbatch_01Rv5hCaNTtUt4JHx1yqoHPL), 2026-11-12; not yet checked by eye.
\scanpage{134-2/41}
\mspage{\hand{26}}

\inserted{en cohomologie étale}
à coéfficients discrets\hand{(}, dans le cas d'une variété projective disons et
de toute section hyperplane, ou d'une variété quasi-projective et de presque
toute section hyperplane (pour ne mentionner que le cas global le plus sim-%
ple), sous les hypothèses de profondeur\hand{(}\inserted[typed]{cohomologique} ``les plus naturelles'' (en fait,
essenti\over{\unsure{a}}{e}llement des conditions néc. et suff. de validité dudit théorème).
Dans le cas commutatif, les techniques de dualité nous suggèrent très clai-%
rement quels sont les meilleurs énoncés possibles, cf. l'exposé de Mme Ray-%
naus dans SGA 2. Mais ces techniques ne valent qu'en se restreignant à des
coéfficients premiers aux caractéristiques, alors que des démonstrations
directes plus géométriques (développées dans SGA 2 av\over{\unsure{e}}{a}nt le développement
du formalisme de la cohomologie étale) donnaient des résultats très voisins
\margin{$\pi_0,\pi_1$}pour le H$^0$ et le H$^1$ (ou le\inserted{$\pi_0$ et le} $\hand{\pi}_1$ , si on préfère) sans telles restrictions,
du moins dans le cas propre (i.e. projectif, au lieu de quasi-projectif).
En fai\hand{t}, ce sont \deleted{\unsure{aussi}} \over{\unsure{d}}{l}es ``résultats les meilleurs possibles'' eux-%
mêmes, énoncés comme conjectures dans SGA 2 dans l'exposé cité de Mme \over{\unsure{R}}{R}ay-%
naud, qui sont démontrés ultérieurement par elle dans sa thèse. Ce qui est
remarquable de notre point de vue, c'est que les énoncés les plus forts se
présentent le plus naturellement sous forme d'énoncés sur des\inserted{1-}\uline{champs} sur
le site étale de la variété algébrique considérée - la notion de ``profon-%
deur $\geq$ i'' (pour i=\over{\unsure{1}}{1},2,3) s'énonçant aussi le plus naturellement en termes
de champs. Non seulement cela, mais alors même qu'on voudrait ignorer la
notion technique de champ et travailler exclusivement en termes de H$^0$ et
H$^1$en utilisant à bloc le formalisme\inserted{cohomologique non commutatif} de Giraud, pour démontrer disons un
théorème de bijectivité $\hand{\pi}_1$(Y)$\hand{\to}\hand{\pi}_1$(X) (ce qui est le résultat le plus pro-%
fond établi dans la thèse de Mme Raynaud), il semble bien qu'on n'y arrive\hand{\unsure{''}}
\xout{xxxx} pas, faute \xout{dxxxxixxa} à ce formalisme d'avoir la souplesse nécessaire\hand{.}
En fait, il faut utiliser comme ingrédients techniques, de façon essentielle\hand{,}
les trois théorèmes suivants directement pour les 1-champs ``de torsion''
(i.e. où les faisceaux en groupes d'automorphismes sont de ind-torsion):
a) théorème de changement de base pour un morphisme propre, b) théorème de
changement de base par un morphisme lisse  c) théorème de \xout{xhanx} ``propreté
cohomologique générique'' pour un morphisme de type fini f: X$\to$S, S intègre
(disant que l'on peut trouver dans S un ouvert U$\neq\emptyset$ tel que pour \uline{tout} change
ment de base S'$\to$ S se factorisant par u, la formule de changement de base
est vraie). \hand{(}Pour b) et c), il faut faire des hypothèses que les fais\unsure{d}eaux
d'automorphismes sont premiers aux cara\unsure{ct}érisitques, et dans c) une hypothè
se de constructtibilité sur ces faisceaux - mais b) et c) ne servent que
dans la version ``générique'' du th. de Lefschetz.\hand{)} C'est avec en vue de telle\hand{s}
applications que Giraud a pris la peine dans son bouquin (si je ne me trompe\hand{)}
de demontrer a)\hand{,} b) (et c) ?) dans le contexte des 1-champs et de leurs

\margin{41}
